The matching statistics of string $q$ with respect to string $r$ is
the array of lengths of longest substring in $q$ that match somewhere
in $r$ \cite{cha94:sub}. Consider for example $q=\ty{AAAATATTAA}$ and
$r=\ty{TTAAAATTAA}$; then the first matching statistic, $m_{\rm
  s}[1]$, is the length of the longest prefix of $q[1...]$ that
matches somewhere in $r$. This happens to be \ty{AAAAT}, which matches
at $r[3]$, so $m_{\rm s}[1]=5$. The next match statistic, $m_{\rm
  s}[2]=4$, and so on, as shown in Table~\ref{tab:ms}.

\begin{table}
  \caption{Matching statistics, $m_{\rm s}$ and match lengths, $m_{\rm
      l}$ of the query $q=\ty{AAAATATTAA}$ with respect to the
    reference $r=\ty{TTAAAATTAA}$.}\label{tab:ms}
  \begin{center}
    \begin{tabular}{ccc}
  \hline
  \# & $m_{\rm s}$ & $m_{\rm l}$\\\hline
1 & 5 & 5\\
2 & 4 & 4\\
3 & 3 & 3\\
4 & 2 & 2\\
5 & 2 & 1\\
6 & 5 & 0\\
7 & 4 & 4\\
8 & 3 & 3\\
9 & 2 & 2\\
10 & 1 & 1\\
\hline
\end{tabular}

  \end{center}
\end{table}

In our program Fur, we implemented something similar to the matching
statistics, which we call the match lengths, $m_{\rm l}$. Instead of
matches starting at every position in $q$, we skip each match before
matching again. Then we interpolate the missing match lengths as
$m_{\rm l}[i]\leftarrow m_{\rm l}[i-1]-1$ whereever $m_{\rm l}[i] <
m_{\rm l}[i-1]-1$.

The first match length is always the same as the first matching
statistics, $m_{\rm s}[1]=m_{\rm l}[1]$. Then the matching length
computation skips forward to start matching again at $q[l+2]$ to find
$m_{\rm l}[l+2]$. Algorithm~\ref{alg:ml} summarizes this.

\begin{algorithm}
  \caption{Calculate match lengths}\label{alg:ml}
  \begin{algorithmic}
    \Require $q$ \Comment{query}
\Require $r$ \Comment{reference}
\Ensure $\ml$ \Comment{array of match lengths along $q$}
\State{$\mathrm{esa}\leftarrow\mathrm{getEsa}(r)$}
\State{$i\leftarrow 1$}
\While{$i\le|q|$} \Comment{matching phase}
   \State{$\ell\leftarrow\mathrm{esa.MatchPrefix}(q[i...])$}
   \State{$\ml[i]\leftarrow\ell$}
   \State{$i\leftarrow i+\ell + 1$}
\EndWhile
\For{$i\leftarrow 2, |q|$} \Comment{interpolation phase}
   \If{$\ml[i]<\ml[i-1]-1$}
      \State{$\ml[i]\leftarrow\ml[i-1]-1$}
   \EndIf
\EndFor

  \end{algorithmic}
\end{algorithm}
